\section{November 11: Rings, Revisited}

\begin{defbox}
    A \bold{ring} is a structure $(R, +, \cdot)$, where $R$ is a set and $+$ and $\cdot$ are binary relations on $R$ such that:
    \begin{enumerate}
        \item $(R, +)$ is an abelian group, and we denote its identity by $0$;
        \item $\cdot$ is associative and has a two-sided identity, denoted $1$; and
        \item $\cdot$ distributes over $+$ from both sides, i.e. for all $a$, $b$, $c \in R$,
        \[ a \cdot (b + c) = a \cdot b + a \cdot c, (a + b) \cdot c = a \cdot c + b \cdot c. \] 
    \end{enumerate}
    Furthermore, we say a ring is \bold{commutative} if $\cdot$ is commutative, i.e. for all $a$, $b \in R$, $a \cdot b = b \cdot a$.
\end{defbox}

Although there is a unique, one element ring where $0 = 1$, we will not consider this ring (call it \textit{trivial}) and in a nontrivial ring, we have $0 \neq 1$.

\begin{defbox}
    An element $u \in R$ is called a \bold{unit} if it is $\cdot$-invertible, i.e. there exists $v \in R$ such that $uv = 1 = vu$. In this case, $v$ is unique and we call it $u^{-1}$. We denote the set of units by $R^\times$. 
\end{defbox}

\begin{defbox}
    An \bold{interval domain} is a nontrivial commutative ring without zero-divisors, i.e. for all $a$, $b \in R$, $a \cdot b = 0 \implies a = 0$ or $b = 0$. A \bold{division ring} is a ring where each nonzero element is a unit, i.e. $R^\times = R \setminus \{0\}$. A \bold{field} is a nontrivial commutative division ring. 
\end{defbox}

\begin{remark}
    By definition, every field is an integral domain. 
\end{remark}

\begin{example}
    Let $R$ be a ring. The ring of $n \times n$ matrices, denoted $M_n(R)$, is the set of all finite sequences of elements of $R$ indexed by pairs $(i,j)$ for $1 \leq i, j \leq n$ and we write them as:
    \[ \begin{pmatrix}
        a_{11} & a_{12} & \dots & a_{1n}\\
        a_{21} & a_{22} & \dots & a_{2n}\\
        \vdots & \vdots & \ddots & \vdots \\
        a_{n1} & a_{n2} & \dots & a_{nn}\\
    \end{pmatrix} \] 
    Where matrix addition is done coordinate-wise, and matrix multiplication is define by the dot product on each cell (these are the same matrix operations we studied in MATH133). 

    For our example, we will typically consider rings of matrices over a field $F$. In this case, there is a nice theory (called linear algebra) which characterises the invertible matrices (units of $M_n(F)$):
    \[ M_n(F)^\times = \text{GL}_n(F) := \{ A \in M_n(F) : \det(A) \neq 0\} \]
    Note that $M_n(F)$ is never commutative for $n \geq 2$. 
\end{example}

\begin{example}
    Let $R$ be a ring. A \textit{formal polynomial over} $R$ is an expression of the form 
    \[ f = a_0 + a_1x + \dots a_nx^n \]
    where $a_i \in R$, $i \in [0, n]$ and $x$ is indeterminate. We say two polynomials are equal if the coefficients $a_0 \dots a_n$ are equal, treating the missing coefficients as zeros. The set of all polynomials over $R$ is denoted by $R[x]$ and it is a ring under $+$ and $\cdot$ operations. 

    The \textit{zero polynomial} is the one with all coefficients zero and we denote it b $0$, and the multiplicative identity is $f = 1$. 
\end{example}